\section{Pendahuluan}

Invoice merupakan dokumen penting dalam proses administrasi penjualan karena berfungsi sebagai dasar penagihan, pencatatan transaksi, dan proses pembayaran pelanggan. Ketepatan pengiriman invoice memengaruhi kelancaran arus kas perusahaan, terutama bagi pelanggan yang menerapkan aturan administrasi tertentu, seperti batas \textit{cut-off}, jadwal penerimaan dokumen, dan prosedur verifikasi sebelum pembayaran dapat diproses. Oleh karena itu, perusahaan membutuhkan sistem yang tidak hanya mampu melacak status pengiriman invoice, tetapi juga memastikan setiap invoice dikirim sesuai dengan ketentuan operasional pelanggan.

Pada perusahaan yang menjadi objek penelitian, proses pelacakan invoice masih dilakukan secara semi-manual menggunakan \textit{spreadsheet}. Staf administrasi harus memeriksa status pengiriman secara manual serta menentukan prioritas pengiriman berdasarkan pengalaman kerja. Kondisi tersebut menyebabkan proses pencarian status invoice menjadi kurang efisien, keputusan prioritas berpotensi berbeda antarstaf, dan pengetahuan operasional sulit ditransfer kepada staf baru. Selain itu, bukti penerimaan (\textit{Proof of Delivery}/POD) belum terdokumentasi secara terintegrasi sehingga proses penelusuran riwayat pengiriman masih membutuhkan waktu yang relatif lama.

Permasalahan tersebut tidak hanya berkaitan dengan pelacakan dokumen, tetapi juga dengan proses pengambilan keputusan. Prioritas pengiriman invoice dipengaruhi oleh berbagai aturan operasional pelanggan, seperti jenis \textit{cut-off}, jadwal penerimaan dokumen, serta ketentuan administrasi lainnya. Sebagian besar pertimbangan tersebut masih tersimpan sebagai \textit{tacit knowledge} yang dimiliki oleh staf administrasi sehingga belum direpresentasikan secara sistematis ke dalam sistem informasi. Akibatnya, proses penentuan prioritas sangat bergantung pada pengalaman individu dan sulit dievaluasi secara konsisten.

Berbagai penelitian telah mengembangkan sistem \textit{document tracking}, \textit{Proof of Delivery} (POD), maupun model klasifikasi berbasis \textit{Decision Tree} pada berbagai domain operasional \cite{10.1007/978-3-642-55032-4_33,gemma,gunadhya_rajawali_logistik_2021,Siahaan_Doni_2025,firmansyah,afifahdedi}. Meskipun demikian, sebagian besar penelitian masih berfokus pada pelacakan status dokumen atau pembangunan model klasifikasi secara terpisah. Penelitian yang mengkaji proses transformasi pengetahuan operasional menjadi dasar pembentukan model keputusan, serta membandingkan pendekatan \textit{knowledge-driven} dan \textit{data-driven} dalam konteks pengelolaan invoice, masih relatif terbatas.

Berdasarkan permasalahan tersebut, penelitian ini mengusulkan sistem pelacakan invoice dan \textit{Proof of Delivery} (POD) berbasis web yang mengintegrasikan proses \textit{Knowledge Acquisition} dan \textit{Knowledge Formalization} untuk merepresentasikan pengetahuan operasional perusahaan. Pengetahuan tersebut kemudian digunakan sebagai dasar penyusunan pedoman \textit{expert labeling} sehingga diperoleh \textit{ground truth} yang digunakan dalam pembangunan dan evaluasi model. Penelitian ini membandingkan dua pendekatan yang berbeda, yaitu model \textit{Rule-Based} sebagai representasi pengetahuan pakar dan algoritma Decision Tree C4.5 sebagai pendekatan berbasis data.

Kontribusi utama penelitian ini meliputi tiga aspek. Pertama, pengembangan sistem pelacakan invoice dan \textit{Proof of Delivery} (POD) berbasis web yang mendukung pencatatan status dan riwayat pengiriman invoice. Kedua, formalisasi pengetahuan operasional menjadi aturan eksplisit yang digunakan sebagai dasar proses \textit{expert labeling}. Ketiga, evaluasi komparatif antara pendekatan \textit{Rule-Based} dan Decision Tree C4.5 untuk mengetahui kemampuan kedua pendekatan dalam merepresentasikan proses pengambilan keputusan pada penentuan prioritas pengiriman invoice.

